%#Split: 05_abilities_daihyo % extract 
%#PieceName: p05_abilities_daihyo
\input{pieces/p05_abilities_daihyo_00}
\section{研究者調書（研究代表者）}
%　　　　＜＜最大　2ページ＞＞

% s14_abilities_S
% s14_pub_S_commands
% begin 研究者情報　＝＝＝＝＝＝＝＝＝＝＝＝＝＝＝＝＝＝＝＝＝＝＝＝＝＝
\renewcommand{\研究者氏名}{市川裕大}
\renewcommand{\研究者氏名ふりがな}{いちかわ ゆうだい}
\renewcommand{\研究者生年月日の年}{1986}
\renewcommand{\研究者生年月日の月}{5}
\renewcommand{\研究者生年月日の日}{31}
\renewcommand{\研究者年齢}{39}
\renewcommand{\研究者所属機関部局職}{\small{東北大学・大学院理学研究科・准教授}}	% use \tiny if necessary
\renewcommand{\研究者学位}{博士（理学）}
% end 研究者情報　＝＝＝＝＝＝＝＝＝＝＝＝＝＝＝＝＝＝＝＝＝＝＝＝＝＝

\input{pieces/s14_pub_S_picture} % pieces

% \PapersInstructions	% <-- 留意事項。これは消すか、コメントアウトしてください。
% \EnvironmentInstructions	% <-- 留意事項。これは消すか、コメントアウトしてください。

%begin 基盤S研究代表者の研究遂行能力及び研究環境 ＝＝＝＝＝＝＝＝＝＝＝＝＝＝＝＝＝＝＝＝
	\noindent
	\textbf{（１）研究代表者のこれまでの研究活動}\\
	私は、2021年に本課題で実施する $\Sigma N$カスプ測定実験（J-PARC E90実験）を課題責任者としてJ-PARCに提案し、シミュレーションを通じた実験立案、検出器の開発など、研究全体を主導してきた。本実験は、私が博士論文で担当した J-PARC E27実験において、$\Sigma N$カスプを副次的に観測したことを契機として構想されたものである（論文\ref{pub:ptep2014}）。また、私はJ-PARCにおいて、様々なハドロン実験を実施しており、筆頭著者として複数の論文を発表しているなど（論文\ref{pub:ptep2024}, \ref{pub:ptep2020}, \ref{pub:ptep2015}）、十分な実績を積んでいる。

	E90実験では、主検出器としてS-2SスペクトロメータとHypTPCを用いる。S-2Sスペクトロメータは、2022年より本格的に建設が始まった新型装置であり、私は現場を取りまとめる責任者の一人としてその建設に大きく貢献した。

	HypTPCは、私が中心となって開発を進めてきた高計数率・大立体角を特徴とする汎用TPC検出器であり、試作機段階から開発の中核を担ってきた。2018年には、放射線医学総合研究所（HIMAC）において陽子ビームを用いた高ビーム強度耐性試験を実験責任者として実施し、HypTPCが想定していた $10^6\,\mathrm{Hz}$ のビーム環境下でも正常に動作することを実証した（論文\ref{pub:nim2019}）。また、$\pi$中間子ビームを想定し、ビーム強度を一桁高めた $10^7\,\mathrm{Hz}$ 環境での対応を目指した改良開発も、科研費若手研究B（16K17714、2016--2017）の支援を受けて並行して実施した。

	さらに、HypTPCを用いた初の物理実験であるHダイバリオン探索実験（E42）を、2021年6月に課題責任者（Co-Spokesperson）として遂行した。現在に至るまで、私はHypTPCを用いた各種実験において、代表者または分担者として中心的な役割を果たしている（例：基盤研究B（代表）21H01097、2021--2025；新学術公募研究（代表）21H00130、2021--2022；基盤研究A（分担）21H04478、2021--2024；基盤研究A（分担）18H03706、2018--2023）。

	また、本研究では、J-PARCにおける $pA$衝突を用いた $\alpha\Lambda$フェムトスコピーを、$\phi$中間子の媒質効果を調べるE88実験のバイプロダクト測定として実施する計画である。E88実験の課題責任者である佐甲と連携し、すでに同実験のコラボレーションに参加しており、フェムトスコピーを主眼とした実験提案書をJ-PARCに共同で提出することを決定している。
	\\

	\noindent
	\textbf{（２）研究代表者の研究環境}\\
	本研究を遂行するにあたっては、HypTPCおよびS-2Sスペクトロメータをはじめとする各種検出器の開発・運用環境が不可欠である。HypTPCは、私が中心となって開発を進めてきた検出器であり、現在もその改良作業を主導している。また、S-2SスペクトロメータはJ-PARC構内に設置されており、共同利用者としてのアクセスが可能であることから、必要な開発環境・設備はすでに整っている。
	\\

	\noindent
	\textbf{（３）研究組織全体の研究環境}\\
	本研究は、J-PARC、LHC（ALICE）、SPring-8、RHIC（STAR）で得られる実験成果を基盤として実施する。J-PARCにおいては、E90実験（市川）およびE88実験（佐甲）がそれぞれ課題責任者のもとで既に進行しており、実験実施に必要な環境は十分に整備されている。SPring-8での $\Lambda p$散乱実験についても、三輪が課題責任者として研究を推進しており、順調に実験が進められている。また、将来の $\Lambda p$散乱実験に向けた検出器開発についても、東北大学やJ-PARC構内の既存実験施設を活用することで、必要な環境は確保されており、研究の遂行に支障はない。

	ALICE実験に関連するデータ解析に関しては、研究分担者の神谷がすでにALICEグループと連携し、フェムトスコピーに関する複数の論文を発表しており、これまでの活動を継続することで、本研究を十分に進めることが可能である。ALICEにおけるフェムトスコピーの研究は主にドイツのミュンヘン工科大学（TUM）が主導しており、研究代表者の市川はすでにTUMの研究者と、本研究の中心課題である $\Sigma N$カスプ分光（E90実験）と $\Sigma p$フェムトスコピーとの相補性について議論を行い、共同研究を進めることで合意している（\url{https://indico.cern.ch/event/1439855/contributions/6461439/}）。

	また、研究分担者の新井田は、これまでRHIC（STAR）実験においても研究を進めており、すでにフェムトスコピーに関する実績を有している。したがって、本研究においても十分な解析体制が整っており、実施に問題はない。
%end 基盤S研究代表者の研究遂行能力及び研究環境 ＝＝＝＝＝＝＝＝＝＝＝＝＝＝＝＝＝＝＝＝

%begin 研究業績リスト ＝＝＝＝＝＝＝＝＝＝＝＝＝＝＝＝＝＝＝＝
	\begin{enumerate}
		\paper{Inclusive spectrum of the d($\pi^+$, $K^+$) reaction at 1.69~GeV/c}{Y. Ichikawa \etal}{PTEP}{2014}{101D03}{2014}
			\label{pub:ptep2014}
				
		\paper{Missing-Mass Measurement of the $^{12}\mathrm{C}(K^-, K^+)$ Reaction at 1.8~GeV/c with the Superconducting Kaon Spectrometer}{Y. Ichikawa \etal}{PTEP}{2024}{091D01}{2024}
			\label{pub:ptep2024}
				
		\paper{An event excess observed in the deeply bound region of the $^{12}\mathrm{C}(K^-, p)$ missing-mass spectrum}{Y. Ichikawa \etal}{PTEP}{2020}{123D01}{2020}
			\label{pub:ptep2020}
		
		\paper{Observation of the ``$K^-pp$''-like structure in the d($\pi^+$, $K^+$) reaction at 1.69~GeV/c}{Y. Ichikawa \etal}{PTEP}{2015}{021D01}{2015}
			\label{pub:ptep2015}

		\paper{High-rate performance of a time projection chamber for an H-dibaryon search experiment at J-PARC}{S.H. Kim, Y. Ichikawa, \etal}{Nucl. Instrum. Meth. A}{940}{359}{2019}
			\label{pub:nim2019}
	\end{enumerate}
%end 研究業績リスト ＝＝＝＝＝＝＝＝＝＝＝＝＝＝＝＝＝＝＝＝

\input{pieces/p05_abilities_daihyo_01}
